% TOP_PROBLEM: {"id": 2416, "title": "Erdős Problem #970", "category_id": 1, "category": {"id": 1, "name": "number_theory", "display_name": "Number Theory", "description": "Properties of integers, prime numbers, Diophantine equations.", "slug": "number-theory", "order_index": 1, "created_at": "2026-07-31T15:26:25.670Z"}, "status": "open"}
% FIRST_POSTED: null
% ENTRY_KIND: statement_only
% Imported from: unsolvedmath_additional_500.tex; UnsolvedMath ID 2416
% !TeX program = lualatex
% Compile twice: lualatex 2416.tex
% Self-contained: no external bibliography, images, or \input files.
% Numeric section labels and visible numbers are original UnsolvedMath IDs.
\documentclass[11pt,a4paper]{article}
\usepackage[margin=24mm,headheight=14pt]{geometry}
\usepackage{fontspec}
\setmainfont{Latin Modern Roman}
\setsansfont{Latin Modern Sans}
\setmonofont{Latin Modern Mono}
\usepackage{amsmath,amssymb,mathtools}
\usepackage{unicode-math}
\setmathfont{Latin Modern Math}
\usepackage{microtype}
\usepackage{enumitem,xurl,needspace}
\usepackage[dvipsnames]{xcolor}
\usepackage{hyperref}
\hypersetup{unicode=true,colorlinks=true,linkcolor=MidnightBlue,urlcolor=MidnightBlue,
 pdftitle={UnsolvedMath Problem 2416},pdfauthor={Source-annotated compilation},bookmarksnumbered=true}
\usepackage{bookmark,tocloft,titlesec,fancyhdr}
\setlength{\cftsecnumwidth}{4.2em}
\cftsetpnumwidth{2.5em}
\cftsetrmarg{4em}
\titleformat{\section}{\Large\bfseries\raggedright}{\thesection}{0.7em}{}
\pagestyle{fancy}\fancyhf{}
\fancyhead[L]{\small\sffamily Additional UnsolvedMath statements}
\fancyhead[R]{\small Original catalogue IDs}
\fancyfoot[C]{\thepage}
\setlength{\parindent}{0pt}
\setlength{\parskip}{5pt plus 1pt}
\setlength{\emergencystretch}{3em}
\setcounter{tocdepth}{1}\setcounter{secnumdepth}{1}
\setlist[itemize]{leftmargin=3em,itemsep=4pt,topsep=4pt,parsep=0pt}
\allowdisplaybreaks
\newcommand{\N}{\mathbb{N}}
\newcommand{\Z}{\mathbb{Z}}
\newcommand{\Q}{\mathbb{Q}}
\newcommand{\R}{\mathbb{R}}
\newcommand{\C}{\mathbb{C}}
\newcommand{\F}{\mathbb{F}}
\newcommand{\A}{\mathbb{A}}
\newcommand{\PP}{\mathbb{P}}
\newcommand{\E}{\mathbb{E}}
\newcommand{\eps}{\varepsilon}
\newcommand{\dd}{\,\mathrm d}
\newcommand{\sref}[2]{\hyperlink{source:#1:#2}{[S#2]}}
\newif\ifshowcatalogue
\showcataloguetrue
\newcommand{\cataloguescope}[1]{\ifshowcatalogue
 \paragraph{Statement recorded in the supplied source.}
 #1\fi}
\begin{document}
% UnsolvedMath ID 2416; source catalogue code EP-970

\setcounter{section}{2415}

\section{A Quadratic Bound for Jacobsthal's Function}\label{2416}

{\small\sffamily UnsolvedMath ID 2416 · Number Theory}

\subsection{Definitions and mathematical statement}

\paragraph{Shared notation.}Unless a section says otherwise, $\N=\{1,2,\ldots\}$,
$\Z$ is the ring of integers, $\Q,\R,\C$ are the rational, real, and complex
fields, $[N]=\{1,\ldots,\lfloor N\rfloor\}$, and $\log$ is the natural
logarithm. For real functions of a parameter tending to infinity,
$f\ll g$ and $f=O(g)$ mean $|f|\leq Cg$ eventually for some fixed $C$;
$f\gg g$ reverses this comparison; $f=o(g)$ means $f/g\to0$;
$f\sim g$ means $f/g\to1$; and $f\asymp g$ means both order bounds.
A subscript on $O$ or $\ll$ specifies allowed constant dependence.
The expression $n^{o(1)}$ in an upper bound means growth smaller than
$n^\epsilon$ for each fixed $\epsilon>0$. Local definitions govern any
exception, finite-field convention, or use of the same symbol for another
object. An asymptotic claim is not automatically an effective algorithm.



Count distinct prime factors of the modulus: $\omega(n)\leq k$. The maximum is over all such positive moduli, including their possible prime powers. Coprimality depends only on the underlying prime set, not those exponents. \sref{2416}{1} \sref{2416}{2}

\paragraph{Statement recorded in the supplied source.}
Let $h(k)$ be Jacobsthal's function, defined to as the minimal $m$ such that, if $n$ has at most $k$ prime factors, then in any set of $m$ consecutive integers there exists an integer coprime to $n$. Determine the order of magnitude of $h(k)$. In particular, is it true that $ h(k) \ll k^2? $ \sref{2416}{1}

\paragraph{Residual task reported by the archive.}

The embedded review (2026-08-17) records: Remove the logarithmic loss from the upper bound or substantially improve the covering lower bound. \sref{2416}{1}

{\small\emph{Transcription note.} Only unambiguous escape/formatting damage and nonmathematical serialization debris were repaired in the displayed source statement.}

\subsection{Short English statement}

How long an interval suffices to contain an integer coprime to every modulus with at most a prescribed number of distinct prime factors, and is a quadratic bound enough?

\subsection{Sources}

{\small

\begin{itemize}

\item[\hypertarget{source:2416:1}{S1}] ULAM.AI, \emph{UnsolvedMath}, supplied \texttt{problems.json}, record 2416 (EP-970), statement and background. The embedded literature-review date is 2026-08-17; that is the archive’s review date, not a fresh certification. Dataset landing page: \url{https://huggingface.co/datasets/ulamai/UnsolvedMath}.

\item[\hypertarget{source:2416:2}{S2}] T. F. Bloom, \emph{Erdős Problems}, Problem 970, \url{https://www.erdosproblems.com/970}. Reference imported from the supplied record; the live problem page was not independently checked for this compilation.

\item[\hypertarget{source:2416:3}{S3}] H. Iwaniec, On the problem of Jacobsthal, Demonstratio Math. 11 (1978), 225-231, \nolinkurl{doi:10.1515/dema-1978-0121.} \url{https://doi.org/10.1515/dema-1978-0121}. Bibliographic information imported from the supplied record.

\item[\hypertarget{source:2416:4}{S4}] K. Ford, B. Green, S. Konyagin, J. Maynard, and T. Tao, Long gaps between primes, J. Amer. Math. Soc. 31 (2018), 65-105, \nolinkurl{doi:10.1090/jams/876.} \url{https://arxiv.org/abs/1412.5029}. Bibliographic information imported from the supplied record.

\item[\hypertarget{source:2416:5}{S5}] Ford, Kevin and Green, Ben and Konyagin, Sergei and Maynard, James and Tao, Terence, Long gaps between primes. J. Amer. Math. Soc. (2018), 65-105. Bibliographic information imported from the supplied record.

\item[\hypertarget{source:2416:6}{S6}] Iwaniec, Henryk, On the problem of {J}acobsthal. Demonstratio Math. (1978), 225--231. Bibliographic information imported from the supplied record.

\end{itemize}

}

\label{end:2416}
\end{document}
